\documentclass[12pt, a4paper]{article}

%=============== PAQUETES ===============
\usepackage[utf8]{inputenc}
\usepackage[T1]{fontenc}
\usepackage[spanish]{babel}
\usepackage{geometry}
\usepackage{amsmath}
\usepackage{graphicx}
\usepackage{xcolor}
\usepackage{hyperref}
\usepackage{tikz}
\usepackage{circuitikz}
\usetikzlibrary{arrows.meta, patterns, babel}


%=============== CONFIGURACIÓN ===============
\geometry{a4paper, margin=2.5cm}

\hypersetup{
    colorlinks=true,
    linkcolor=blue,
    urlcolor=cyan,
    pdftitle={Análisis de Examen de Electrónica},
}

\definecolor{darkgreen}{rgb}{0.0, 0.5, 0.0}
\definecolor{darkred}{rgb}{0.8, 0.0, 0.0}

\newcommand{\acierto}[1]{\textcolor{darkgreen}{\textbf{Acierto:} #1}}
\newcommand{\error}[1]{\textcolor{darkred}{\textbf{Error:} #1}}
\newcommand{\faltante}[1]{\textcolor{blue!70!black}{\textbf{Faltó añadir:} #1}}

%=============== DOCUMENTO ===============

\begin{document}

\title{Análisis de Examen de Dispositivos Electrónicos}
\author{Evaluación de Jesús López}
\date{11 de Noviembre de 2025}
\maketitle

\section*{Análisis General}
El estudiante demuestra una comprensión inicial de algunos conceptos básicos de los diodos, pero presenta errores conceptuales significativos en la mayoría de las preguntas. Se observan confusiones en las características I-V, los diagramas de circuitos, las formas de onda de salida y las definiciones de los modelos teóricos. Es necesario un repaso profundo de los fundamentos.

\hrulefill
\vspace{1cm}

\section{Pregunta 1: Diferencias entre diodo real e ideal}

\subsection*{Respuesta del Estudiante (Resumen)}
El estudiante menciona que el diodo real tiene una caída de tensión que no existe en el ideal. Afirma incorrectamente que la característica I-V de ambos es una curva exponencial.

\subsection*{Evaluación}
\begin{itemize}
    \item \acierto{La mención de la caída de tensión en el diodo real (aprox. 0.7V) frente a la ausencia de esta en el ideal (0V) es correcta y fundamental.}
    \item \error{La característica I-V del diodo ideal \textbf{no es exponencial}. Es una función escalón: se comporta como un interruptor perfecto.}
    \item \faltante{No se mencionan otras diferencias clave como la corriente de fuga en polarización inversa, la tensión de ruptura (breakdown) o la resistencia interna del diodo real.}
\end{itemize}

\subsection*{Respuesta Correcta}
Las diferencias fundamentales son:
\begin{enumerate}
    \item \textbf{Caída de Tensión ($V_f$):} El diodo ideal tiene $V_f = 0$V. El diodo real tiene una caída de tensión (aprox. 0.7V para silicio).
    \item \textbf{Polarización Inversa:} El diodo ideal no conduce corriente ($I_R = 0$). El diodo real tiene una pequeña corriente de fuga ($I_S$).
    \item \textbf{Tensión de Ruptura ($V_{BR}$):} El diodo ideal soporta cualquier tensión inversa. El diodo real tiene una tensión de ruptura a partir de la cual conduce abruptamente en inversa.
    \item \textbf{Característica I-V:} La del ideal es una función escalón. La del real es una curva exponencial descrita por la ecuación de Shockley.
\end{enumerate}

\begin{figure}[h!]
    \centering
    \begin{tikzpicture}[scale=0.8, xshift=-0.5cm]
        % Ejes
        \draw[->] (-3,0) -- (3,0) node[right] {$V_D$};
        \draw[->] (0,-1) -- (0,4) node[above] {$I_D$};
        
        % Diodo Ideal (rojo)
        \draw[thick, red] (0,0) -- (0,3.5);
        \draw[thick, red] (-2.5,0) -- (0,0);
        \node[red, below right] at (1,2) {Ideal};

        % Diodo Real (azul)
        \draw[thick, blue, domain=0:0.8, smooth] plot (\x, {0.05*(exp(8*\x) - 1)});
        \draw[thick, blue] (-2.5, -0.05) -- (0, -0.05);
        \node[blue, above right] at (1, 2.5) {Real};
        \draw[dashed] (0.7,0) -- (0.7, 1.5);
        \node[below] at (0.7,0) {0.7V};
    \end{tikzpicture}
    \caption{Comparación de las curvas I-V de un diodo real y uno ideal.}
\end{figure}

\newpage

\section{Pregunta 2: Rectificador de media onda}

\subsection*{Respuesta del Estudiante (Resumen)}
Explica que el circuito no deja pasar los semiciclos negativos. Dibuja un circuito rectificador con filtro de capacitor y una forma de onda de salida correspondiente a un rectificador de \textbf{onda completa}.

\subsection*{Evaluación}
\begin{itemize}
    \item \acierto{La explicación de que el diodo bloquea un semiciclo es correcta.}
    \item \error{El diagrama del circuito es incorrecto. Se pidió un rectificador de media onda simple, no uno con filtro.}
    \item \error{La forma de onda de salida dibujada es incorrecta. Corresponde a un rectificador de onda completa, que aprovecha ambos semiciclos.}
\end{itemize}

\subsection*{Respuesta Correcta}
El rectificador de media onda simple consta de una fuente AC, un diodo y una resistencia de carga ($R_L$). Durante el semiciclo positivo de la señal de entrada, el diodo se polariza en directa y permite el paso de la corriente. Durante el semiciclo negativo, se polariza en inversa y bloquea la corriente.

\begin{figure}[h!]
    \centering
    \begin{circuitikz}[scale=0.9, transform shape]
        \draw (0,0) to[sV, l=$V_{in}$] (0,2) to[short] (1,2) to[D, l=D] (3,2) to[R, l=$R_L$] (5,2) to[short] (6,2) node[ground]{};
        \draw (0,0) to[short] (6,0) to[short] (6,2);
    \end{circuitikz}
    \caption{Diagrama de un rectificador de media onda simple.}
\end{figure}

\begin{figure}[h!]
    \centering
    \begin{tikzpicture}[scale=1.2]
        % Señal de entrada
        \draw[->] (-0.2,0) -- (4,0) node[below] {$t$};
        \draw[->] (0,-1.2) -- (0,1.2) node[left] {$V$};
        \draw[blue, thick] (0,0) sin (1,1) cos (2,0) sin (3,-1) cos (4,0);
        \node[above] at (2,1.2) {Entrada ($V_{in}$)};

        % Señal de salida
        \begin{scope}[yshift=-2.5cm]
            \draw[->] (-0.2,0) -- (4,0) node[below] {$t$};
            \draw[->] (0,-1.2) -- (0,1.2) node[left] {$V_{out}$};
            \draw[red, thick] (0,0) sin (1,1) cos (2,0);
            \draw[red, thick, dashed] (2,0) sin (3,-1) cos (4,0); % Semiciclo negativo recortado
            \draw[red, thick] (2,0) -- (4,0);
            \node[above] at (2,1.2) {Salida ($V_{out}$)};
        \end{scope}
    \end{tikzpicture}
    \caption{Formas de onda de entrada y salida para un rectificador de media onda.}
\end{figure}

\newpage

\section{Pregunta 4: Ventaja del rectificador de onda completa}

\subsection*{Respuesta del Estudiante (Resumen)}
Afirma que la ventaja es que el rectificador de onda completa 'mantiene la señal de salida constante'.

\subsection*{Evaluación}
\begin{itemize}
    \item \error{La afirmación es incorrecta. La salida de un rectificador de onda completa (sin filtro) no es constante, sino una señal DC pulsante.}
\end{itemize}

\subsection*{Respuesta Correcta}
La principal ventaja del rectificador de onda completa sobre el de media onda es su \textbf{mayor eficiencia}:
\begin{enumerate}
    \item \textbf{Mayor Tensión DC de Salida:} Aprovecha ambos semiciclos de la señal AC, resultando en un valor promedio (DC) en la salida que es el doble que el de un rectificador de media onda.
    \item \textbf{Filtrado más Sencillo:} La frecuencia del rizado (ripple) en la salida es el doble de la frecuencia de entrada. Una frecuencia de rizado más alta es mucho más fácil y barata de filtrar para obtener una señal DC suave.
\end{enumerate}

\section{Pregunta 5: Modelo de tramos lineales}

\subsection*{Respuesta del Estudiante (Resumen)}
Describe un método para señales con 'saltos' o discontinuidades y dibuja una onda escalonada, lo cual no tiene relación con el modelo del diodo.

\subsection*{Evaluación}
\begin{itemize}
    \item \error{La respuesta es completamente incorrecta. Confunde el modelo del dispositivo con un tipo de señal.}
\end{itemize}

\subsection*{Respuesta Correcta}
El \textbf{modelo de tramos lineales} (o lineal por partes, \textit{piecewise linear model}) es una \textbf{aproximación} de la curva exponencial del diodo real utilizando segmentos de línea recta para simplificar el análisis de circuitos.

El modelo más común aproxima el diodo como:
\begin{itemize}
    \item Un \textbf{interruptor abierto} si el voltaje a través de él es menor que la tensión umbral ($V_D < V_{th}$, ej. 0.7V).
    \item Una \textbf{fuente de tensión ideal de $V_{th}$} en serie con una resistencia ($r_D$) si el voltaje intenta superar $V_{th}$. A menudo, para una primera aproximación, se considera $r_D = 0$.
\end{itemize}
Este modelo es más preciso que el modelo ideal (que asume $V_{th}=0$) pero mucho más simple de analizar matemáticamente que el modelo exponencial completo, siendo muy útil para cálculos manuales.

\begin{figure}[h!]
    \centering
    \begin{tikzpicture}[scale=0.8]
        % Ejes
        \draw[->] (-1,0) -- (3,0) node[right] {$V_D$};
        \draw[->] (0,-0.5) -- (0,4) node[above] {$I_D$};
        
        % Modelo lineal por tramos
        \draw[thick, blue] (0.7,0) -- (2, 2.6); % Resistencia r_D > 0
        \draw[thick, blue] (-0.5,0) -- (0.7,0);
        \node[blue, right] at (1.5, 3) {Modelo Lineal por Tramos};
        \draw[dashed] (0.7,0) -- (0.7, 2);
        \node[below] at (0.7,0) {$V_{th}$};
    \end{tikzpicture}
    \caption{Característica I-V del modelo de tramos lineales.}
\end{figure}

\end{document}